\section{Related Work}\label{sec:related}

\subsection{Contrastive Vision-Language Pre-training}
Recent advances in foundational vision-language models have fundamentally transformed multimedia representation learning. The paradigm was largely established by CLIP \cite{radford2021learning}, which demonstrated the efficacy of contrastive learning on web-scale image-text pairs. Subsequent iterations, such as OpenCLIP \cite{ilharco2021openclip} trained on LAION-5B \cite{schuhmann2022laion}, scaled these approaches further. More recently, SigLIP \cite{zhai2023sigmoid} introduced a pairwise sigmoid loss, mitigating the batch size dependency inherent to the standard InfoNCE loss. Concurrently, Multimodal Large Language Models (MLLMs) like Qwen-VL \cite{bai2023qwen} have integrated vision encoders with large causal language models to achieve unprecedented visual reasoning capabilities. Despite these successes, mapping high-dimensional multimodal signals into a single, monolithic embedding space invariably introduces representation blind spots, particularly regarding fine-grained object attributes and subtle semantic nuances. In contrast, Vecna proposes a heterogeneous multi-encoder ensemble, leveraging orthogonal architectural inductive biases to construct a more robust, composite feature space.

\subsection{Dense, Sparse, and Hybrid Information Retrieval}
The retrieval literature has long been bifurcated into sparse, lexical approaches and dense, semantic paradigms. Sparse models like BM25 \cite{robertson2009probabilistic} excel at exact keyword matching but fail to capture latent semantics. Conversely, dense retrieval paradigms, initialized by models like SBERT \cite{reimers2019sentence}, project text into dense semantic spaces, enabling robust synonymy resolution. Modern architectures such as BGE-M3 \cite{chen2024bge} unify these modalities, offering multi-granularity embeddings that simultaneously encode dense semantic structures, sparse lexical features, and multi-vector representations. To fuse disparate retrieval signals, rank aggregation techniques like Reciprocal Rank Fusion (RRF) \cite{cormack2009reciprocal} are frequently employed. However, Vecna eschews heuristic late-fusion rank aggregation in favor of a mathematically rigorous, unified convex combination of visual cosine similarity, dense BGE-M3 semantic scoring, and sparse BM25 lexical alignment, further augmented by exact-phrase boundary verification to eliminate semantic drift during fusion.

\subsection{Generative Query Expansion and Hypothetical Document Embeddings}
Query expansion techniques have historically relied on pseudo-relevance feedback to augment query representations. With the advent of large language models, paradigms such as Hypothetical Document Embeddings (HyDE) \cite{gao2023precise} have emerged. HyDE employs an LLM to generate a hypothetical, albeit potentially hallucinated, document in response to a zero-shot query, subsequently utilizing the dense embedding of this synthetic document for retrieval. While effective in bridging the semantic gap, such unconstrained generative expansion is acutely susceptible to hallucination, introducing spurious features that degrade retrieval precision in factual domains. Vecna addresses this limitation by introducing the Strict Factual Boundary principle. By constraining the generative expansion process strictly to observable visual attributes and verifiable metadata, Vecna ensures that the synthesized pseudo-documents remain anchored in empirical reality, thereby mitigating the hallucination problem inherent to unconstrained HyDE approaches.

\subsection{Temporal Sequence Retrieval in Video}
Video retrieval introduces the orthogonal challenge of temporal dynamics. End-to-end architectures, such as Frozen in Time \cite{bain2021frozen}, employ computationally expensive 3D Convolutional Neural Networks or spatio-temporal Transformers to jointly model spatial features and temporal evolution. While these models capture complex intra-video temporal dependencies, their quadratic complexity relative to sequence length renders them computationally prohibitive for large-scale video corpora. Vecna bypasses the latency bottlenecks of end-to-end temporal cross-attention by adopting a decoupled architecture. We propose a lightweight, frame-level dynamic programming algorithm incorporating a two-pointer sliding window mechanism. This formulation allows Vecna to efficiently identify temporally contiguous semantic segments, achieving temporal localization capabilities comparable to heavy spatio-temporal transformers but at a fraction of the computational cost.
